% !TEX root = task-3-11-1.tex
% Editable task bank. Append new BankTask blocks at the end of the document.
\documentclass[12pt,a4paper]{article}
\usepackage{physics-textbook}
\usepackage{graphicx,needspace}
\usetikzlibrary{decorations.pathmorphing}
\newif\ifShowSolutions
\ifdefined\TutorEdition
  \ShowSolutionstrue
\else
  \ShowSolutionsfalse % Change to \ShowSolutionstrue to display solutions.
\fi
\newcounter{banktask}
\newsavebox{\BankTaskBox}
\newcommand{\BankTaskContent}[3]{%
  \noindent\textbf{\thebanktask.}\quad #2\par
  \ifShowSolutions
    \Needspace{50mm}% Keep the solution heading with its opening text or diagram.
    #3\par
    \ifdefined\TutorEdition
      \ifcsname TutorSource#1\endcsname
        \nopagebreak\smallskip{\small\csname TutorSource#1\endcsname\par}%
      \fi
    \fi
  \fi
}
\newcommand{\BankTask}[3]{%
  \refstepcounter{banktask}\par\medskip
  \sbox{\BankTaskBox}{\begin{minipage}{\linewidth}%
    \BankTaskContent{#1}{#2}{#3}\end{minipage}}%
  \ifdim\dimexpr\ht\BankTaskBox+\dp\BankTaskBox\relax>\dimexpr\textheight-20mm\relax
    \Needspace{35mm}\BankTaskContent{#1}{#2}{#3}%
  \else
    \Needspace{\dimexpr\ht\BankTaskBox+\dp\BankTaskBox\relax}%
    \noindent\usebox{\BankTaskBox}\par
  \fi
  \medskip
}
\begin{document}
\BookSubsection{3.11.1}{Работа тока}


% Source sdamgia-1406: https://phys-ege.sdamgia.ru/problem?id=1406
\BankTask{sdamgia-1406}{
Чему равно время прохождения тока силой $5$ А по проводнику, если при напряжении на его концах $120$ В в проводнике выделяется количество теплоты, равное $540$ кДж? (Ответ дайте в секундах.)
}{
\smallskip\noindent\textbf{Решение}\par
Согласно закону Джоуля — Ленца количество выделяемого тепла в проводнике определяется выражением $Q=UIt.$ Отсюда находим время прохождения тока по проводнику:

\[t= \frac{Q}{UI} = \frac{540\,\text{\,\text{кДж}}}{120\,\text{В} \cdot  5\,\text{А}} = \frac{540 \cdot  10 ^3\,\text{Дж}}{120\,\text{В} \cdot  5\,\text{А}} =900\,\text{с}.\]

Ответ: $900$.
}

% Source sdamgia-4910: https://phys-ege.sdamgia.ru/problem?id=4910
\BankTask{sdamgia-4910}{
Комната освещается четырьмя одинаковыми параллельно включенными лампочками. Расход электроэнергии за час равен Q. Каким должно быть число параллельно включенных лампочек, чтобы расход электроэнергии в час был равен $2$Q?
}{
\smallskip\noindent\textbf{Решение}\par
При параллельном соединении ко всем лампочкам приложено одинаковое напряжение $U.$ Независимо от общего количества лампочек за $1$ час каждая лампочка расходует одинаковое количество энергии $\frac{U ^2}{R} \cdot  1\,\text{ч},$ где $R$ — сопротивление лампочки. Согласно условию, расход $4$ лампочек равен $Q=4 \cdot  \frac{U ^2}{R} \cdot  1\,\text{ч}.$ Следовательно, чтобы расход энергии за час стал равен $2Q,$ необходимо, чтобы было включено $8$ лампочек (параллельно соединенных).

Ответ: $8$.
}

% Source sdamgia-19720: https://phys-ege.sdamgia.ru/problem?id=19720
\BankTask{sdamgia-19720}{
Показания квартирного электросчетчика за некоторое время возросли на $200$ кВт · ч. Какую работу совершил за это время электрический ток, протекавший через установленные в данной квартире электроприборы? Ответ запишите в мегаджоулях.
}{
\smallskip\noindent\textbf{Решение}\par
Работа электрического тока равна потребленной электроэнергии, которая по условию равна $200$ кВт · ч. 

Переводим эту величину в Джоули: $200$ кВт · ч = $200$ · $1000\,\text{Вт}$ · $3600\,\text{с} = 720$ · $10$\textasciicircum{}\{$6$\} Дж = $720$ МДж.

Ответ: $720$.
}

% Source sdamgia-42639: https://phys-ege.sdamgia.ru/problem?id=42639
\BankTask{sdamgia-42639}{
При прохождении по проводнику постоянный электрический ток силой $3$ А в течение $60\,\text{с}$ совершает работу, равную $270$ кДж. Чему равно сопротивление проводника? Ответ запишите в омах.
}{
\smallskip\noindent\textbf{Решение}\par
Из формулы работы электрического тока $A=I ^2 Rt$ находим сопротивление проводника:

\[R= \frac{A}{I ^2 t} = \frac{270000}{9 \cdot  60} =500\,\text{Ом}.\]

Ответ: $500$.
}

% Source sdamgia-46771: https://phys-ege.sdamgia.ru/problem?id=46771
\BankTask{sdamgia-46771}{
Паяльник подключили к источнику постоянного напряжения $220$ В. Какое количество теплоты выделится за $20$ минут работы паяльника, если сила тока в цепи его питания равна $1{,}2$ А? Ответ дайте в килоджоулях.
}{
\smallskip\noindent\textbf{Решение}\par
Количество теплоты, выделяющееся в проводнике при прохождении тока равно работе тока

\[Q = UIt = 220 \cdot  1{,}2 \cdot  1200 = 316800\,\text{Дж} = 316{,}8\,\text{\,\text{кДж}}.\]

Ответ: $316{,}8$.
}


% Source 0990BD: https://ege.fipi.ru/bank/questions.php?search=1&proj=BA1F39653304A5B041B656915DC36B38&qid=0990BD
\BankTask{0990BD}{
При прохождении по проводнику постоянный электрический ток силой $3\,\text{А}$ в течение $60\,\text{с}$ совершает работу, равную $270\,\text{кДж}$. Чему равно сопротивление проводника (в $\text{Ом}$)?
}{
\smallskip\noindent\textbf{Решение от GPT}\par
Работа постоянного тока $A=I^2Rt$. Следовательно,
\[
R=\frac{A}{I^2t}=\frac{270\cdot10^3}{3^2\cdot60}=500\,\text{Ом}.
\]
Ответ: $500\,\text{Ом}$.
}

\end{document}
